% STARTER DRAFT for Conor -- written with Claude Code on 2026-10-20 at Patrick's
% request. Not included in main.tex. Use, rewrite or ignore; see drafts/README.md.

% ---------------------------------------------------------------------
% -> sections/related_work.tex (shared with Emmett; Emmett's notes are in emmett.tex)
% Sources: readings/jacovi_goldberg_2020.md, lanham_2023.md, arcuschin_2025.md,
% pramanik_2026_nsf_cot.md; arXiv abstracts for turpin2023, chen2025, lyu2023,
% parcalabescu2024, xiong2025 (NO READING NOTES YET for these five -- read before citing).
% ---------------------------------------------------------------------
\textbf{Faithfulness.} Jacovi and Goldberg separate faithfulness, which asks whether
an explanation reflects the model's actual reasoning, from plausibility, which asks
whether it convinces a reader. They argue faithfulness should be graded rather than
binary~\cite{jacovi2020}.

\textbf{Testing chain-of-thought faithfulness.} Lanham et al.\ intervene on the
reasoning, by truncating it, adding mistakes or paraphrasing it, and find that models
vary widely in how much their answers depend on it~\cite{lanham2023}. Turpin et al.\
show that biasing features added to the prompt change answers without being mentioned
in the explanation~\cite{turpin2023}. Chen et al.\ find that reasoning models reveal
their use of prompt hints in at least 1\% of cases where they use them, but often in
fewer than 20\%~\cite{chen2025}. Arcuschin et al.\ show that unfaithful reasoning also
occurs on natural, non-adversarial prompts~\cite{arcuschin2025}. Xiong et al.\ perturb
the thinking drafts of reasoning models and find that final answers often fail to
align with the draft's conclusion~\cite{xiong2025}.
% CHECK: confirm xiong2025 is the paper the group meant (see refs.bib).

\textbf{Formal and symbolic checking.} Faithful CoT translates a question into a
symbolic program and runs it with a deterministic solver, so the reasoning chain
explains the answer by construction~\cite{lyu2023}. NSF-CoT checks each reasoning step
for groundedness, validity and utility, using SMT-style verification combined with
LLM-based entailment~\cite{pramanik2026}. Unlike both, we leave the model's own
natural-language working unchanged and check it.

\textbf{Consistency versus faithfulness.} Parcalabescu and Frank argue that most
faithfulness tests measure self-consistency at the output level, not faithfulness to
the model's inner workings~\cite{parcalabescu2024}. We return to this in
Section~\ref{sec:discussion}.

% ---------------------------------------------------------------------
% -> sections/method.tex, subsection "Logic checker and FOLIO validation"
% Sources: README.md "FOLIO logic checker"; journal/2026-10-16.md; journal/2026-10-20.md
% ---------------------------------------------------------------------
For first-order logic we use Z3~\cite{demoura2008}. The checker parses premises and a
conclusion written in first-order logic (quantifiers, predicates and relations,
equality, and the usual connectives) and reports whether the conclusion is entailed,
contradicted or unknown. Inconsistent premises, parse errors and solver timeouts are
reported separately from logical uncertainty, so checker failures are never counted
as model errors. Optional formal reasoning steps are checked in order, and only proved
steps can support later ones. We validate the checker on the official FOLIO
annotations~\cite{han2024folio} (Table~\ref{tab:folio}). This is a check of the
checker, not a faithfulness result: we did not check any model's logic reasoning,
because free-text reasoning would first have to be translated into first-order logic.

% ---------------------------------------------------------------------
% -> sections/results.tex, subsection "Logic checker validation on FOLIO"
% Source: outputs/clean/folio_validation_summary.json
% ---------------------------------------------------------------------
With a 500\,ms limit per solver call, 199 of the 204 validation rows parsed, and the
checker's verdict matched the dataset label for 191 of them (Table~\ref{tab:folio}).
Five rows had malformed annotations, and no call timed out or found inconsistent
premises. The eight disagreements have not yet been audited against the
natural-language text.

% ---------------------------------------------------------------------
% -> sections/conclusion.tex
% Sources: outputs/clean/paper_stats.json; journal/2026-10-07-circuit-tracing.md;
% journal/2026-10-16.md; 50-problem report "Data integrity"
% ---------------------------------------------------------------------
\textbf{Limitations.} We test one model on one dataset, with 50 problems per mode, so
the intervals are wide. The checker sees only calculations written as equations or
with operator words. A step such as ``she has 9 left'' is invisible to it, and 17 of
the 100 responses could not be compared at all. Its coverage, estimated from hand
counts, is \textbf{TBD}. Steps within a response are not independent, so the
step-level interval is optimistic. The planted-working and copy-vs-use tests use
Qwen3-0.6B, while the GSM8K outputs come from 1.7B. Circuit-tracing graphs describe a
transcoder replacement of the model, not Qwen3 itself. Finally, the FOLIO results
validate the logic checker only.

\textbf{Future work.} Natural errors are too rare on GSM8K to measure unfaithfulness
directly. Harder problems, where working and answers diverge more often, would give
the test more power. Translating free-text reasoning into first-order logic, with
human review, would let us check Qwen3 on FOLIO. Feature interventions on attribution
graphs could test whether the answer depends on the working.
% Add the hint test / copy-vs-use here if either missed the Wed 21 Oct 9pm freeze.

\textbf{Conclusion.} On GSM8K, Qwen3's final answers never contradicted its own
written arithmetic, with or without thinking. That shows the model is consistent, not
that it is faithful. Telling the two apart needs tests that change the reasoning or
the prompt, and our results from those are \textbf{TBD}.
